\documentclass[11pt]{article}
\usepackage[margin=1in]{geometry}
\usepackage{amsmath,amssymb,amsthm,booktabs}
\usepackage[hidelinks]{hyperref}
\usepackage{xurl}
\setlength{\emergencystretch}{3em}
\newtheorem{theorem}{Theorem}
\newtheorem{lemma}[theorem]{Lemma}
\newtheorem{corollary}[theorem]{Corollary}
\newcommand{\rank}{\operatorname{rank}}
\title{A Coxeter reflection obstruction for the four-braid knot $K_*$}
\author{Written proof and exact polynomial certificate}
\date{17 November 2026}
\begin{document}
\maketitle

The knot $K_*$ defined below has no meridional Coxeter reflection-class
quotient whose class requires four generators. This includes finite
and infinite Coxeter groups. The obstruction comes from a polynomial
identity in a general symmetric Gram matrix with six free entries.
The Gram matrix may be singular. The result concerns common symmetric
forms in characteristic zero; it gives no augmentation-rank gap or
upper bound for unrestricted KCH representations.

\section{The knot and the reflection statement}
Use the literal four-braid word
\begin{equation}\label{eq:beta}
 \beta=(1,2,-3,-3,2,-3,-3,2,1,1,2,2,2,-3,2,2,2).
\end{equation}
Here a letter $\pm i$ acts on positions $i-1,i$, numbered from zero.
The local PD convention is fixed by the following port table:
\[
\begin{array}{c|cccc}
 &0&1&2&3\\\hline
 +i&\mathrm{in}_L&\mathrm{out}_L&\mathrm{out}_R&\mathrm{in}_R\\
 -i&\mathrm{in}_R&\mathrm{in}_L&\mathrm{out}_L&\mathrm{out}_R .
\end{array}
\]
Opposite ports are joined through a crossing; ports one and three
form the overpass. Close the bottom endpoints to the top endpoints
in their original positions. This is the Spherogram braid-closure
convention. Let $K_*$ be this closure. The strand permutation is a
four-cycle, so the closure is a knot. The retained seventeen-row PD
is verified directly from the port table and endpoint gluings.
No census name or numerical geometric identification is used.

Let $k$ be a characteristic-zero field, $B$ a symmetric bilinear form
on a $k$-vector space, and
\[
             R_u(x)=x-B(u,x)u,\qquad B(u,u)=2 .
\]
Then $R_u$ is an involution preserving $B$, and its displacement image
is $ku$. The form on the ambient space, and on the span of the
normal vectors, may be singular.

\begin{theorem}\label{thm:reflection}
In every representation of $\pi_1(S^3\setminus K_*)$ taking meridians
to nonisotropic orthogonal reflections preserving one symmetric
bilinear form over a characteristic-zero field, the span of all
meridional normal directions has dimension at most three.
The assertion also holds for the mirror of $K_*$.
\end{theorem}

Conjugacy of meridians permits the normal vectors to be chosen as
images of one normal vector. Their common nonzero norm can be made
two by rescaling $B$; this requires no square root in $k$.
The proof will rule out independence of the four top-meridian normals.
Their span is invariant under the four corresponding reflections,
which generate the image of the knot group. All conjugate meridian
normals then lie in that same span. This explains the global span
statement once independence has been excluded.

\section{Root propagation and closing signs}
The positive Hurwitz move and its inverse are
\begin{equation}\label{eq:moves}
 H_i:(u,v)\longmapsto(R_u(v),u),\qquad
 H_i^{-1}:(u,v)\longmapsto(v,R_v(u)).
\end{equation}
For an involutory meridian assignment with incoming values $a,b$,
the positive literal crossing has PD values $(a,b,bab,b)$ and outgoing
values $(b,bab)$. The negative literal crossing has PD values
$(b,a,aba,a)$ and outgoing values $(aba,a)$.
Thus the Hurwitz word for \eqref{eq:beta} is
\begin{equation}\label{eq:w}
 w=(-1,-2,3,3,-2,3,3,-2,-1,-1,-2,-2,-2,3,-2,-2,-2).
\end{equation}
This determines the sign convention by local crossing relations,
rather than by a name for the braid.

Assume the four top normals are independent and use them as a basis
$e_0,e_1,e_2,e_3$ of their span. Their Gram matrix is
\begin{equation}\label{eq:gram}
 G=\begin{pmatrix}
 2&a&b&c\\a&2&d&e\\b&d&2&f\\c&e&f&2
 \end{pmatrix}.
\end{equation}
All six entries $a,b,c,d,e,f$ are unrestricted. In particular no
Gram determinant is assumed nonzero.

For normalized vectors,
\begin{equation}\label{eq:sign}
                R_u=R_v\quad\Longleftrightarrow\quad v=\pm u .
\end{equation}
Equality first forces the nonzero displacement lines to coincide,
and the equal norms then force the proportionality scalar to have
square one. Each root move in \eqref{eq:moves} has determinant $-1$
on the two moving root columns: it swaps them and adds a multiple of
one column to the other. Consequently all intermediate root frames
remain independent. A closing reflection tuple therefore has final
roots $\varepsilon_i e_i$ with
\begin{equation}\label{eq:product}
                    \prod_{i=0}^3\varepsilon_i=-1.
\end{equation}

\begin{lemma}\label{lem:sign-gauge}
For the word $w$, all eight sign patterns satisfying
\eqref{eq:product} are equivalent under changes
$e_i\mapsto\delta_i e_i$, where $\delta_i\in\{1,-1\}$.
On Gram entries these are polynomial coordinate changes
$G_{ij}\mapsto\delta_i\delta_jG_{ij}$.
\end{lemma}
\begin{proof}
The initial strand ending in position $i$ is
$p(i)$, where $(p(0),p(1),p(2),p(3))=(1,3,0,2)$.
The identity
$R_{\delta u}(\eta v)=\eta R_u(v)$
shows that the final root in position $i$ acquires sign $\delta_{p(i)}$.
Relative to the new initial frame, its closing sign is therefore
\[
                \varepsilon'_i
                    =\delta_{p(i)}\delta_i\varepsilon_i .
\]
The permutation $p$ is a four-cycle. Prescribe $\delta$ at one position
and solve these equations successively around the cycle. The final
equation holds exactly when the two sign products agree. Both choices
of the initial global sign work. This proves the claim using only
sign changes; it introduces no denominator.
\end{proof}

\section{The exact polynomial contradiction}
It is enough to take
$\varepsilon=(-1,1,1,1)$.
Write $F_v$ for propagation by the letters of a word $v$, in their
listed order. Split $w$ after its seventh letter:
\[
 A=(-1,-2,3,3,-2,3,3),\qquad
 C=(-2,-1,-1,-2,-2,-2,3,-2,-2,-2).
\]
Let $L=F_A(e_0,e_1,e_2,e_3)$ and
$T=F_{C^{-1}}(-e_0,e_1,e_2,e_3)$, where $C^{-1}$ reverses the
letter order and changes each sign. Because the root moves are
invertible, the signed closure condition is exactly $L=T$.

Compute these two tuples in
$\mathbb Z[a,b,c,d,e,f]^4$ using \eqref{eq:gram}.
List the sixteen coordinate differences of $L-T$ in position order,
and within each vector in coordinate order. Divide each nonzero
difference by its positive integer coefficient content.
The resulting primitive polynomials $P_0,\ldots,P_{15}$ have degrees
\[
 (32,31,30,24,44,43,42,36,36,37,38,33,25,26,27,22).
\]
Only nonzero integer constants have been divided out; this is valid
in every characteristic-zero field.

\begin{lemma}\label{lem:identity}
The retained integer polynomials $Q_0,\ldots,Q_{15}$ satisfy
\begin{equation}\label{eq:identity}
                   \sum_{j=0}^{15}Q_jP_j=16
             \quad\hbox{in }\mathbb Z[a,b,c,d,e,f].
\end{equation}
They have a total of $75\,136$ nonzero monomials.
\end{lemma}
\begin{proof}
The exact finite certificate stores rational multipliers $q_j$
whose denominators divide $16$. Put $Q_j=16q_j$; these have integer
coefficients. Every coefficient and exponent is retained.
The standard-library checker reconstructs $L,T$ from the literal
word, regenerates the primitive $P_j$, and compares all of their
coefficients with the retained records. It then multiplies and sums
the integer sparse polynomials and obtains the constant polynomial
$16$. The monomial exponents are nonnegative and there are only the
six variables in \eqref{eq:gram}. Thus this is a polynomial identity,
with no determinant localization or saturation. The full coefficient
replay is the finite computational part of the proof.
\end{proof}

Every signed closure in the chosen branch would make all $P_j$
zero, contradicting \eqref{eq:identity}, since $16\ne0$ in characteristic
zero. Lemma~\ref{lem:sign-gauge}
covers the other seven branches. This excludes four independent
top-meridian normals and proves the span conclusion of
Theorem~\ref{thm:reflection}. A mirror homeomorphism induces an
isomorphism of knot groups carrying meridians to meridians up to
inversion. Inversion does not change reflection images. The span of
all meridional normal directions is consequently unchanged, proving
the mirror assertion as well.

\section{The Coxeter consequence}
Here a reflection in a Coxeter group means a conjugate of a simple
generator. For a conjugacy class $\mathcal C$ generating a group,
its class-generating rank is the minimum number of its members
generating that group.

\begin{corollary}\label{cor:coxeter}
If a knot group of $K_*$, or of its mirror, surjects onto a Coxeter
group and takes a meridian to a reflection, then the image meridian
class has generating rank at most three. This applies to finite and
infinite Coxeter groups. In particular such a quotient cannot supply
a class-generating-rank lower bound four.
\end{corollary}
\begin{proof}
For completeness the geometric representation needed here follows
directly from the Coxeter presentation. Give the simple-root basis
$\{\alpha_s\}$ the form
\[
 B(\alpha_s,\alpha_s)=2,\qquad
 B(\alpha_s,\alpha_t)=
 \begin{cases}
 -2\cos(\pi/m_{st}),&m_{st}<\infty,\\
 -2,&m_{st}=\infty.
 \end{cases}
\]
Its normalized simple-root reflections square to the identity.
For each finite $m_{st}$ the two-root plane is nondegenerate and
has an orthogonal complement. The product on the plane has
eigenvalues $\exp(\pm2\pi i/m_{st})$; for $m_{st}=2$ it is $-I$.
It fixes the orthogonal complement. Thus all defining relations
hold and give a representation, even if the full form is singular.
There is no pair-product relation when $m_{st}=\infty$.
Faithfulness is not needed below.

Let $U$ be the span of the four top-meridian normal directions in
this representation. Theorem~\ref{thm:reflection} gives $\dim U\le3$.
Each of these four reflections acts trivially on the quotient by
$U$, so every element of the surjective Coxeter image does.
For a simple reflection $s$,
$(R_{\alpha_s}-I)V=k\alpha_s$; hence every $\alpha_s$ lies in $U$.
The simple-root basis therefore has at most three members.

The abelianization of a Coxeter group is
$(\mathbb Z/2)^{o}$, where $o$ is the number of components of the
graph joining simple generators with odd finite Coxeter label.
This follows immediately by abelianizing its presentation:
an odd label equates two generators and an even label gives no
additional relation. A group generated by one conjugacy class of
involutions has cyclic abelianization, so $o=1$.
Simple generators joined by an odd-labelled edge are conjugate.
Consequently all simple generators lie in the image meridian class.
They generate the group in at most three members of that class.
\end{proof}

\begin{corollary}\label{cor:sums}
For every integer $s\ge0$, put $K_s=K_*\mathbin{\#}{}^s3_1$.
In a characteristic-zero representation of its knot group by
nonisotropic reflections preserving one symmetric bilinear form,
the span of all meridional normal directions has dimension at most
$3+s$. Every surjective Coxeter reflection-class quotient of this
knot group has class-generating rank at most $3+s$.
Thus it cannot supply a meridional lower bound $4+s$.
\end{corollary}
\begin{proof}
Restrict the representation to the $K_*$ factor.
Theorem~\ref{thm:reflection} applies in the same ambient space, so
its meridional normal directions span at most three dimensions.
The standard three-crossing trefoil diagram has the Wirtinger
presentation
\[
                  \langle x,y\mid xyx=yxy\rangle ,
\]
with $x$ and $y$ both meridians. A connected-sum decomposing annulus
identifies a meridian of each factor. Van Kampen's theorem therefore
gives the knot group as a product amalgamated along that meridian.
For each trefoil choose $x$ as the identified meridian; only the
normal of $y$ can add a direction beyond those already present.
The images are generated by the $K_*$ meridians and these $s$
additional reflections. Their normal span is invariant under all
of them, so it contains every conjugate meridional normal direction.
This gives the bound $3+s$.

For a full Coxeter quotient, apply its geometric representation.
Every simple-root direction lies in this normal span by the
displacement argument in Corollary~\ref{cor:coxeter}; the Coxeter
rank is therefore at most $3+s$. Its odd-labelled graph is connected
by the same cyclic-abelianization argument. Its simple generators
belong to the image meridian class, proving the stated class bound.
This proof uses no meridional-rank equality or cover calculation for
the connected sums.
\end{proof}

\section{Replay, controls and scope}
The package consists of the literal braid/PD input,
the rational unit certificate, a standard-library checker and its
receipt. The default replay checks the braid port construction
against all $68$ retained PD darts, its sphere embedding and one
component, the four actual top-meridian arcs, and a nonconstant
involutory colouring of all seventeen crossings. It verifies all
eight sign transports and the complete polynomial identity.

A genuine positive control is supplied by $T(4,5)$.
The four star roots $e_0-e_j$ in $\mathbb Q^5$ have Gram matrix with
diagonal two and all offdiagonal entries one; their reflections
are the four star transpositions of $S_5$.
For the Hurwitz word $(1,2,3)^5$ they close with signs $(-1,-1,-1,1)$.
The opposite word has a control using the four adjacent-difference
roots, with Gram matrix of type $A_4$.
Both controls have independent root frames. The checker verifies
their root and permutation trajectories, inverse moves, and all
eight sign-gauged Gram specializations for each control.
An altered multiplier coefficient, a mismatched braid sign,
a corrupted PD edge and a negative monomial exponent are rejected
by default.

Sage/Singular was used to discover and extract the identity;
neither is required for its replay. The written sign-gauge,
normal-direction and Coxeter arguments remain essential mathematical
parts. The receipt pins this source, the input, certificate and
checker. Dated internal PDF-review records are preserved under replay
with their own source and PDF hashes; a replay does not transfer an
old review to a new edition.

This companion makes no census identification, crossing-minimality,
geometric-family or worldwide-priority assertion. Those questions
belong to the separate knot study. The common symmetric-form
hypothesis is essential here: no absence of unrestricted KCH
representations, augmentation-rank upper bound, or augmentation gap
is claimed. The result is an exact internal computation with written
proofs, not formal verification or external specialist review.
\end{document}
